\documentclass[10pt]{article}

\usepackage[utf8]{inputenc}

\usepackage[T1]{fontenc}

\usepackage{amsmath}

\usepackage{amsfonts}

\usepackage{amssymb}

\usepackage[version=4]{mhchem}

\usepackage{stmaryrd}

\usepackage{graphicx}

\usepackage[export]{adjustbox}

\graphicspath{ {./images/} }



\begin{document}

н о в Н. В., Таблицы математической статистики, 2 изд., м.

$1968 ;$;

С м и р н о в Н. В.. "Докл. АН СССР", 1941, т. 33 , № 5 , ОІІИБОЧНОГО ДЕКОДИРОВАНИЯ ВЕРОЯТНОСТЬ - одна из возможных мер характеризации точности воспроизведения сообений, передаваемых по каналу связи (см. также Сообщений точность воспроизведения). Пусть для передачи сообщения $\xi$, вырабатываемого источником сообений и принимающего $M$ различных возможных значений $1,$. . ., $M$ с распределением вероятностей $p_{m}=\mathrm{P}\{\xi=m\}, \quad m=1, \ldots ., M_{1}$ используется нек-рый канал связи. Тогда для фиксированных методов кодирования и декодирования (см. также Инбормачии передача) О. Д. в. $P_{e, m}, m=1, \cdots, M$, определяется равенством



\[

P_{e, m}=\mathrm{P}\{\tilde{\xi} \neq m \mid \xi=m\}

\]



где $\tilde{\xi}$ - декодированное сообщение на выходе канала.



$\mathrm{C} \mathrm{p}$ е д н е й $О$. д. в. $P_{e}$ наз. величина



\[

P_{e}=P\{\tilde{\xi} \neq \xi\}=\sum_{m=1}^{M} p_{m} P_{e, m} .

\]



Особый интерес представляет изучение о п т и м а л ьн о й О. Д. в. $P_{e}^{\text {opt }}$, определяемой равенством



\[

P_{\mathrm{e}}^{\mathrm{opt}}=\inf P_{e},

\]



где нижняя грань берется по всевозможным методам кодирования и декодирования. Ввиду трудностей получения точного выражения для $P_{e}^{\mathrm{opt}}$, связанных с тем обстоятельством, что в общем случае неизвестны оптимальные методы кодирования, интерес представляет изучение асимптотич. поведения $P_{e}^{\mathrm{opt}}$ при возрастании длительности передачи по каналу. Точнее, рассматривается следующая ситуация. Пусть для передачи используется отрезок длины $N$ канала связи с дискретным временем и $R=(\ln M) / N$. Требуется изучить асимптотич. поведение $P_{e}^{\mathrm{opt}}$ при $N \rightarrow \infty$ и $R=$ const (это означает, что длительность Іередачи возрастает, а ее скорость остается постоянной). Если рассматриваемый отрезок канала является отрезком однородного канала без паляти с дискретным времевем и конечными пространствами $Y$ п $\tilde{Y}$ значений компонент сигналов на входе и выходе, то известны следующие верхние и нижние оценки $P_{e}^{\text {opt }}=P_{e}^{\mathrm{opt}}(N, R)$ :



\[

\begin{gathered}

\exp \left\{-N\left[E_{s p}(R-\alpha(N))+\beta(N)\right]\right\} \leqslant \\

\leqslant P_{\mathrm{e}}^{\mathrm{opt}}(N, R) \leqslant \exp \left\{-N E_{r}(R)\right\},

\end{gathered}

\]



где $\alpha(N)$ и $\beta(N)$ стремятся к нулю с ростом $N$, а функции $E_{r}(R)$ и $E_{s p}(R)$ определяются следующим образом:



\[

\begin{gathered}

E_{r}(R)=\max _{0} \max _{q}\left[E_{0}(\rho, q(\cdot))-\rho R\right], \\

E_{s p}(R)=\sup _{\rho>0} \max _{q(\cdot)}\left[E_{0}(\rho, q(\cdot))-\rho R\right],

\end{gathered}

\]



где



$E_{0}(\rho, q(\cdot))=-\ln \sum_{\tilde{y} \in \bar{Y}}\left[\sum_{y \in Y} q(y) q(y, \bar{y})^{\frac{1}{1+\rho}}\right]^{1+\rho} \cdot$



Здесь $q(\cdot)=\{q(y), y \in Y\}-$ произвольное распределение вероятностей на $Y, q(y, \tilde{y})=\mathbf{P}\left\{\tilde{\eta}_{k}=\tilde{y} \mid \eta_{k}=y\right\}$, $y \in Y, \quad \tilde{y} \in \tilde{Y}$, а $\eta_{k}$ и $\tilde{\eta}_{k}, k=1, \ldots, N$, - компоненты сигналов на входе и выходе рассматриваемого канала без памяти. Известно, что $E_{r}(R)$ и $E_{s p}(R)$, определяемые формулами (2) и (3), являются положительными, выпуклыми вниз, монотонно убывающими функциями



\includegraphics[max width=\textwidth, center]{2023_07_08_823b3a0c15137b4c9135g-1}

способность, причем $E_{r}(R)=E_{s p}(R)$ при $R \geqslant R_{c r}$ здесь $R_{c r}$ - величина, определяемая переходной матрицей канала и называемая к р и т и ч е с к о й с к ор о с т ь ю для данного канала без памяти. Таким образом, для значений $R, R_{c r} \leqslant R<C$, главные члены верхней и нижней оценок (1) для $P_{e}^{\mathrm{opt}}$ асимптотически совпадают, откуда следует, что для этого диапазона изменения $R$ известно точное значение ф у н к ц и и на д е ж н о с т и к а н а л а $E(R)$, определяемой равенством



\[

E(R)=\varlimsup_{N \rightarrow \infty} \frac{-\ln P_{\mathrm{e}}^{\mathrm{opt}}(N, R)}{N} .

\]



Для значений $R, 0 \leqslant R<R_{c r}$, точное значение $E(R)$ для произвольных каналов без памяти неизвестно (1983), хотя оценки (1) и могут быть улучшены. Экспоненциальный характер убывания $P_{e}^{\mathrm{opt}}(N, R)$ доказан для весьма широкого класса каналов.



Лum.: [1] Г а л л а $\Gamma$ е p P., Теория информации и надежная связь, пер. с англ., М., $1974 ; ;[2] \Phi$ а й н с т е й н А., Основы



теории информации, пер. о англ., М., 1960.





\end{document}